
\chapter{Report on the trace formula}\label{II}

In this text, I have attempted to present Grothendieck's cohomological theory
of $L$-functions as directly as possible. I follow very closely some of the
lectures given by Grothendieck at the IHES in the spring of 1966. In the
spirit of this volume, no appeal will be made to SGA 5---apart from two
references to passages in Exposé XI, which is otherwise independent of that
seminar.

In \S 10 (pp. 81--90) of \cite{de73}, the reader will find a summary of the
theory in the crucial case of curves.










\section{Some reminders}\label{II:1}





\subsection{Notation}\label{II:1-1}

$p,q,\dF,\dF_q$: $p$ is a prime number, $q=p^f$ a power of $p$, and $\dF$
an algebraic closure of the prime field $\dF_p$. For every power $p^n$ of
$p$, denote by $\dF_{p^n}$ the subfield with $p^n$ elements of $\dF$.

$X_0,X$: An object over $\dF_q$ will often be denoted by a symbol carrying a
subscript $0$. The same symbol without the $0$ then denotes the object
obtained from it by extension of scalars to $\dF$. For example, if $\sF_0$ is
a sheaf on a scheme $X_0$ over $\dF_q$, denote by $\sF$ its inverse image on
$X=X_0\otimes_{\dF_q} \dF$.

$F$: If $X_0$ is a scheme over $\dF_q$, the endomorphism of $X_0$ that
``sends the point with coordinates $x$ to the point with coordinates $x^q$''
is called the \emph{Frobenius morphism} and denoted by $F$: it is the identity
on the underlying topological space and, for $x$ a local section of the
structure sheaf, $F^* x = x^q$.

The endomorphism of $X$ obtained from it by extension of scalars is also
denoted by $F$. On the set $X(\dF) = X_0(\dF)$ of rational points of $X$,
$F$ acts as the Frobenius substitution
$\varphi\in\gal(\dF/\dF_q)$ defined by $\varphi(x)=x^q$. In particular,
$X^F=X_0(\dF_q)$.

In case of ambiguity, we shall write $F_{(q)}$ rather than $F$.










\subsection{Frobenius correspondence}\label{II:1-2}

Let $X_0$ be a scheme over $\dF_q$, and let $\sF_0$ be a sheaf on $X_0$. I
like to view the \emph{Frobenius correspondence} (an $F$-endomorphism of
$\sF_0$) as follows:
\begin{equation*}\tag{1.2.1}\label{II:eq:1-2-1}
  F^* : F^*\sF_0 \iso \sF_0
\end{equation*}

Regard $\sF_0$ as the sheaf of local sections of a space $[\sF_0]$ étale over
$X_0$ ($[\sF_0]$ is an algebraic space in the sense of M. Artin). The
functoriality of $F$ gives a commutative diagram
\[\xymatrix{
  [\sF_0] \ar[r]^-F \ar[d]
    & [\sF_0] \ar[d] \\
  X_0 \ar[r]^F
    & X_0
}\]
Because $[\sF_0]$ is étale over $X_0$, this diagram is Cartesian; it gives an
isomorphism $[\sF_0]\iso F^*[\sF_0]$ whose inverse is
\eqref{II:eq:1-2-1}. We also denote by $F^*$ the correspondences or morphisms
deduced from $F$ by extension of scalars or by functoriality, such as
$F^*:\h_c^\bullet(X,\sF)\to \h_c^\bullet(X,\sF)$.

For a formal definition, I refer to \cite[XI.1,2]{sga5}. Only the case
$q=p$ is considered there, but the general case is treated in the same way.






\subsection{}\label{II:1-3}

The Frobenius correspondence is functorial in $(X_0,\sF_0)$. If
$f_0:X_0\to Y_0$ is a morphism of $\dF_q$-schemes, $\sF_0$ a sheaf on $X_0$,
$\sG_0$ a sheaf on $Y_0$, and
$u\in\hom(f_0^*\sG_0,\sF_0)$ an $f_0$-morphism from $\sG_0$ to $\sF_0$, there
is a commutative diagram of spaces $f_0 F = F f_0$, and above it a
commutative diagram of sheaves
\[\xymatrix{
  X_0 \ar[r]^-{f_0} \ar[d]^-F
    & Y_0 \ar[d]^-F \\
  X_0 \ar[r]^-{f_0}
    & Y_0
}\qquad\qquad
\xymatrix{
  \sF_0
    & \ar[l]_-u \sG_0 \\
  \sF_0 \ar[u]_-{F^*}
    & \ar[l] \sG_0 \ar[u]^-{F^*}
}\]
In particular, the correspondence on $f_{0*}\sF_0$ deduced from the
Frobenius correspondence of $(X_0,\sF_0)$ is the Frobenius correspondence of
$(Y_0,f_{0*}\sF_0)$.






\subsection{}\label{II:1-4}

The natural morphism $Y\to Y_0$ is a limit of étale morphisms. It follows
that for every abelian sheaf $\sF_0$ on $X_0$, the inverse images on $Y$ of
the $\eR^if_0\sF_0$ are the $\eR^if_*\sF$.

On $\eR^if_*\sF$, one has
\begin{enumerate}[\indent a)]
  \item the correspondence
    $F^*\eR^i f_*\sF \iso \eR^i f_*\sF \xrightarrow{F^*} \eR^if_*\sF$
    deduced by functoriality from the Frobenius correspondence;
  \item the correspondence $F^*\eR^if_*\sF\to\eR^if_*\sF$ that is the
    inverse image on $Y$ of the Frobenius correspondence on
    $\eR^i f_{0*}\sF_0$.
\end{enumerate}

It follows from (\ref{II:1-3}) that these correspondences coincide. They
induce the initial term of a Frobenius action on the entire spectral sequence
\[
  \h^p(Y,\eR^q f_*\sF) \rightarrow \h^{p+q}(X,\sF) \text{.}
\]
Likewise for cohomology with proper support, and for the spectral sequence
\[
  \h_c^p(Y,\eR^q f_* \sF) \rightarrow \h_c^{p+q}(X,\sF)
\]





\subsection{Change of finite field}\label{II:1-5}

Let $(X_0,\sF_0)$ be over $\dF_q$, let $n$ be an integer, and let
$(X_1,\sF_1)$ over $\dF_{q^n}$ be obtained from $(X_0,\sF_0)$ by extension of
scalars. On the one hand there is the Frobenius correspondence
$F_{(q)}:X_0\to X_0)$, $F_{(q)}^*:F_{(q)}^*\sF_0\to \sF_0$, and on the other
there is $F_{(q^n)}:X_1\to X_1$,
$F_{(q^n)}^*:F_{(q^n)}^*\sF_1\to\sF_1$. It is easily verified that
$F_{(q^n)}$ is obtained from the $n$-th power of $F_{(q)}$ by extension of
scalars from $\dF_q$ to $\dF_{q^n}$; after extension of scalars to $\dF$, one
still has the same identity $F_{(q^n)}=F_{(q)}^n$ between endomorphisms of
$X$ and correspondences on $X$.

If $x$ is a closed point of $X_0$, with $[k(x),\dF_q]=n$, and
$\bar x\in X_0(\dF) = X(\dF)$ lies over $x$, then $\bar x$ is a fixed point
of $F_{(q^n)}=F_{(q)}^n$. Denote by $F_x^*$ the endomorphism of
$\sF_{\bar x}$ induced by $F_{(q^n)}^*$. Up to isomorphism,
$(F_{\bar x}^*,\sF_{\bar x})$ does not depend on the choice of $\bar x$.
The trace, etc. of $F_x^*$ is denoted by notation such as
$\tr(F_x^*,\sF)$.






\subsection{\texorpdfstring{$L$}{L}-function}\label{II:1-6}

Let $X_0$ be a scheme of finite type over $\dF_q$. Denote by $|X_0|$ the set
of its closed points, and for $x\in |X_0|$ set
$\deg(x)=[k(x),\dF_p]$. If $\sF_0$ is a $\dQ_\ell$-sheaf on $X_0$ (for this
notion, see below)---or a constructible flat sheaf of modules over a
commutative Noetherian ring $A$---denote by $L(X_0,\sF_0)$ the formal series
with constant term $1$, in $\dQ_\ell\llbracket t\rrbracket$ or in
$A\llbracket t\rrbracket$,
\[
  L(X_0,\sF_0) = \prod_{x\in|X_0|} \det\left(1-F_x^* t^{\deg(x)},\sF\right)^{-1} \text{.}
\]






\subsection{Remark}\label{II:1-7}

In definition (\ref{II:1-6}), one has $\deg(x)=[k(x),\dF_p]$ (absolute
degree), and not $\deg(x)=[k(x),\dF_q]$ (degree relative to $\dF_q$). The
effect is that $L(X_0,\sF_0)$ depends only on the scheme $X_0$ and the sheaf
$\sF_0$, not on $q$ and the structure morphism
$X_0\to\spec(\dF_q)$. The indeterminate $t$ is an avatar of $p^{-s}$.






\subsection{The Galois point of view}\label{II:1-8}

This subsection (n$^\circ$) will not be used in the remainder of the report. For $\sF_0$
an abelian sheaf on $X_0$, the Galois group $\gal(\dF/\dF_q)$ acts on
$\dF$, and by transport of structure on $\h_c^i(X,\sF)$. In particular, the
Frobenius substitution $\varphi$ (\ref{II:1-1}) induces an automorphism, also
denoted by $\varphi$, of $\h_c^i(X,\sF)$. The Frobenius correspondence, on the
other hand, defines an endomorphism $F^*$ of $\h_c^i(X,\sF)$. One has
\begin{equation*}\tag{1.8.1}\label{II:eq:1-8-1}
  {F^*}^{-1} = \varphi \qquad \text{(on $\h_c^i(X,\sF)$).}
\end{equation*}

To see this, apply (\ref{II:1-4}) to the case where $Y_0$ is a point:
$Y_0=\spec(\dF_q)$. One finds that $\h_c^i(X,\sF)$ is the fiber, at the
geometric point $\spec(\dF)$ of $Y_0$, of $\eR^i f_! \sF_0$, and that $F^*$
is the fiber of the Frobenius correspondence of $\eR^i f_!\sF_0$. Let
$[\eR^i f_! \sF_0]$ be as in (\ref{II:1-2}). One has
$\h_c^i(X,\sF) = [\eR^i f_!\sF_0](\dF)$ and
\begin{enumerate}[\indent a)]
  \item $\varphi$ acts by transport of structure, via its action on $\dF$;
  \item $F^*$ is the inverse of
    $F:[\eR^i f_!\sF_0](\dF) \to [\eR^i f_!\sF_0](\dF)$.
\end{enumerate}

The conclusion follows from the identity $F=\varphi$ of (\ref{II:1-1}).

One advantage of the Galois point of view is that it permits reasoning by
transport of structure.










\section{\texorpdfstring{$\dQ_\ell$}{Ql}-sheaves}\label{II:2}

The notion of $\dQ_\ell$-sheaves is developed in detail in
\cite[V,VI]{sga5}. We shall use only the definitions and results below.






\begin{definition_}\label{II:2-1}
Let $X$ be a Noetherian scheme. A $\dZ_\ell$-sheaf $\sF$ on $X$ is an
inverse system of sheaves $\sF_n$ ($n\geqslant 0$), with $\sF_n$ a
constructible sheaf of $\dZ/\ell^{n+1}$-modules \cite[IX.2]{sga4}, such that
the transition morphism $\sF_n\to\sF_{n-1}$ factors through an isomorphism
\[
  \sF_n\otimes_{\dZ/\ell^{n+1}} \dZ/\ell^n \iso \sF_{n-1}\text{.}
\]
The sheaf $\sF$ is said to be \emph{lisse} if the $\sF_n$ are locally
constant.
\end{definition_}

The terminology of \cite{sga5} is different: ``constructible
$\dZ_\ell$-sheaf'' is used there for ``$\dZ_\ell$-sheaf,'' and
``constructibly twisted constant'' for ``lisse.''

Likewise below, for $\dQ_\ell$-sheaves.






\subsection{}\label{II:2-2}

If $k$ is a separably closed field, a sheaf of $\dZ/\ell^n$-modules on
$\spec(k)$ is identified with a $\dZ/\ell^n$-module, and the functor
$\sF\mapsto \varprojlim \sF_n$ identifies the $\dZ_\ell$-sheaves on
$\spec(k)$ with the $\dZ_\ell$-modules of finite type. This justifies the
following definition: the \emph{fiber} of the $\dZ_\ell$-sheaf $\sF$ at the
geometric point $x$ of $X$ is the $\dZ_\ell$-module
$\varprojlim (\sF_n)_x$.






\subsection{}\label{II:2-3}

Suppose $X$ is connected, and let $x$ be a geometric point of $X$. It is
known that the ``fiber at $x$'' functor is an equivalence between the
category of constructible locally constant sheaves of
$\dZ/\ell^n$-modules and the category of finite-type
$\dZ/\ell^n$-modules equipped with a continuous action of $\pi_1(X,x)$.
Passing to the limit gives the following proposition.






\begin{proposition_}\label{II:2-4}
Under the hypotheses above, the ``fiber at $x$'' functor is an equivalence
between the category of lisse $\dZ_\ell$-sheaves on $X$ and that of
finite-type $\dZ_\ell$-modules equipped with a continuous action of
$\pi_1(X,x)$.
\end{proposition_}






\begin{proposition_}\label{II:2-5}
Let $\sF$ be a $\dZ_\ell$-sheaf on a Noetherian scheme $X$. There is a finite
partition of $X$ into locally closed parts $X_i$ such that the $\sF|X_i$ are
lisse.
\end{proposition_}

Set $\gr_\ell^n\sF=\ell^n\sF_m/\ell^{m+1}\sF_m$ for $m\geqslant n$, and
$\gr_\ell\sF=\bigoplus \gr_\ell^n\sF$. It is a sheaf of modules over the
Noetherian ring $\gr_\ell\dZ=\dZ/\ell[t]$, the graded ring of $\dZ$ for the
$\ell$-adic filtration. It is a quotient of
$\gr_\ell^0\sF \otimes_{\dZ/\ell} \dZ/\ell[t]$, hence a constructible sheaf
of $\gr_\ell\dZ$-modules \cite[IX.2]{sga4}; and there is a finite partition
of $X$ into locally closed parts $X_i$ such that
$\gr_\ell \sF|X_i$ is locally constant. Each of the sheaves
$\gr_\ell^n\sF|X_i$ is then locally constant, as are the $\sF_m|X_i$, which
are successive extensions of them.






\subsection{}\label{II:2-6}

If $\sF$ and $\sG$ are two $\dZ_\ell$-sheaves on $X$, set
$\hom(\sF,\sG) = \varprojlim(\sF_n,\sG_n)$; this is a $\dZ_\ell$-module.
Since $\hom(\sF_n,\sG_n) = \hom(\sF_m,\sG_n)$ for $m\geqslant n$, one also
has
\[
  \hom(\sF,\sG) = \varprojlim_n \varinjlim_m \hom(\sF_m,\sG_n) \text{,}
\]
so that the functor $\sF\mapsto \text{``}\varprojlim\text{''}\sF_n$ is fully
faithful from the category of $\dZ_\ell$-sheaves to that of pro-sheaves on
$X$ ($=$ pro-objects of the category of sheaves on $X$). We shall use it to
identify $\dZ_\ell$-sheaves with certain pro-sheaves (those such that every
$\sF_n=\sF/\ell^{n+1}\sF$ is a sheaf, and
$\sF\iso \text{``}\varprojlim\text{''}\sF_n$). It is easily verified that if
$\sF+\text{``}\varprojlim\text{''}\sF_n$, with $\sF_n$ a constructible sheaf
of $\dZ/\ell^n$-modules, then $\sF$ is a $\dZ_\ell$-sheaf if and only if the
inverse systems (in $n$) $\sF_n/\ell^{k+1}\sF_n$ are essentially constant:
the inverse system of
$\sF_k'=\varprojlim_n\sF_n/\ell^{k+1}\sF_n$ then satisfies the conditions of
(\ref{II:2-1}), and
$\sF=\text{``}\varprojlim\text{''}\sF_k'$.

In what follows, we abandon the notation $\sF_n$ of (\ref{II:2-1}); if
$\sF$ is a $\dZ_\ell$-sheaf, the sheaf $\sF_n$ of (\ref{II:2-1}) will be
denoted by $\sF\otimes\dZ/\ell^{n+1}$.






\subsection{}\label{II:2-7}

In the abelian category of pro-sheaves, the subcategory of
$\dZ_\ell$-sheaves is stable under kernels, cokernels, and extensions. This is
clear for cokernels. For kernels, if $f:\sF\to\sG$ is a morphism of
$\dZ_\ell$-sheaves, it follows readily from
(\ref{II:2-4}, \ref{II:2-5}) and Artin--Rees that the inverse system of
$\ker(f:\sF\otimes\dZ/\ell^n\to\sG\otimes\dZ/\ell^n)$ satisfies the
criterion given in (\ref{II:2-6}) above. There is even an integer $r$ such
that for every $n$
\[
  \ker(f)\otimes\dZ/\ell^n = \ker(f:\sF\otimes\dZ/\ell^{n+r}\to\sG\otimes\dZ/\ell^{n+r})\otimes\dZ/\ell^n \text{.}
\]
The case of extensions is left to the reader.






\subsection{}\label{II:2-8}

A $\dZ_\ell$-sheaf $\sF$ is \emph{torsion-free} if
$\ker(\ell:\sF\to\sF) = 0$. Each $\sF\otimes\dZ/\ell^n$ is then flat over
$\dZ/\ell^n$ (i.e. its fibers are finite-rank free
$\dZ/\ell^n$-modules). It follows from (\ref{II:2-4}, \ref{II:2-5}) that for
every $\dZ_\ell$-sheaf $\sF$ there is an integer $n$ such that
$\sF/\ker(\ell^n)$ is torsion-free.






\subsection{}\label{II:2-9}

The abelian category of \emph{$\dQ_\ell$-sheaves} on $X$ is the quotient of
that of $\dZ_\ell$-sheaves by the subcategory of torsion
$\dZ_\ell$-sheaves. Equivalently, its objects are the $\dZ_\ell$-sheaves,
and, denoting by $\sF\otimes\dQ_\ell$ the $\dZ_\ell$-sheaf $\sF$, viewed as
an object of the category of $\dQ_\ell$-sheaves, one has
\[
  \hom(\sF\otimes\dQ_\ell,\sG\otimes\dQ_\ell) = \hom(\sF,\sG)\otimes_{\dZ_\ell}\dQ_\ell \text{.}
\]

The \emph{fiber} of a $\dQ_\ell$-sheaf $\sF\otimes\dQ_\ell$ at a geometric
point $x$ of $X$ is the finite-dimensional $\dQ_\ell$-vector space
$\sF_x\otimes_{\dZ_\ell}\dQ_\ell$.






\subsection{}\label{II:2-10}

Let $X$ be a separated scheme of finite type over an algebraically closed
field $k$, and let $\sF$ be a constructible $\dQ_\ell$-sheaf on $X$. Let
$\sF'$ be a constructible $\dZ_\ell$-sheaf such that
$\sF\sim \sF'\otimes_{\dZ_\ell}\dQ_\ell$. We shall see in
(\ref{II:4-11}) that
\[
  \left(\varprojlim \h_c^q(X,\sF'\otimes \dZ/\ell^n)\right)\otimes_{\dZ_\ell}\dQ_\ell
\]
is a finite-dimensional $\dQ_\ell$-vector space. It depends only on $\sF'$,
and is denoted by $\h_c^q(X,\sF)$. The proof does not use the fact that
$\ell$ is prime to the characteristic, but only that case concerns us.






\subsection{}\label{II:2-11}

In this exposé, all important calculations will be carried out at the level
of $\dZ/\ell^n$-sheaves, with passage to $\dQ_\ell$-sheaves taking place only
at the very end. In practice, however, it is often convenient to work
systematically with $\dQ_\ell$-sheaves. In Exposés V and VI of
\cite{sga5}, Jouanolou gives the formalism that makes this possible. One of
the essential results (\hyperref[VI:2-2]{VI 2.2}) is that if $f:X\to Y$ is a separated morphism
of finite type of Noetherian schemes, and $\sF$ a $\dZ_\ell$-sheaf on $Y$,
then, for every $q$, the pro-sheaf
$\text{``}\varprojlim\text{''}\eR^q f_! (\sF\otimes\dZ/\ell^n)$ is a
$\dZ_\ell$-sheaf on $Y$. He even proves a stronger regularity property (of
Artin--Rees type) for the inverse system of the
$\eR^q f_!(\sF\otimes\dZ/\ell^n)$.

Jouanolou even works in a more general setting, encompassing the very useful
case of $E_\lambda$-sheaves, for $E_\lambda$ a finite extension of
$\dQ_\ell$. They are defined as the $\dQ_\ell$-sheaves were, with
$\dQ_\ell$ replaced by $E_\lambda$, $\dZ_\ell$ by the ring $\sO_\lambda$ of
the valuation $\lambda$, and $\dZ/\ell^n$ by $\sO_\lambda/\pi^n$, where
$\pi$ is a uniformizer. It would amount to the same thing to define an
$E_\lambda$-sheaf as a $\dQ_\ell$-sheaf $\sF$ equipped with a homomorphism
$E_\lambda\to\operatorname{End}(\sF)$.









\section{Trace formula and \texorpdfstring{$L$}{L}-functions}\label{II:3}

\emph{We use the notation of \S\ref{II:1}, and $\ell$ is a prime number
$\ne p$.}

Let $X_0$ be a separated scheme of finite type over $\dF_q$ ($q=p^f$), and
let $\sF_0$ be a constructible $\dQ_\ell$-sheaf on $X_0$. Grothendieck's
cohomological interpretation of $L$-functions is the following theorem:






\begin{theorem_}\label{II:3-1}
$L(X_0,\sF_0) = \prod_i \det\left(1-F^* t^f, \h_c^i(X,\sF)\right)^{(-1)^{i+1}}$.
\end{theorem_}

It is obtained as a consequence of the trace formula:






\begin{theorem_}\label{II:3-2}
For every integer $n$,
$\sum_{x\in X^{F^n}} \tr\left({F^n}^*,\sF_x\right) = \sum_i (-1)^i \tr\left({F^n}^*,\h_c^i(X,\sF)\right)$.
\end{theorem_}






Let us briefly recall the classical method for deducing (\ref{II:3-1}) from
(\ref{II:3-2}). Set $T=t^f$; both sides are formal series in $T$ with
constant term $1$. Since $\dQ_\ell$ has characteristic $0$, it suffices to
prove that the logarithmic derivatives $T\frac{d}{dT}\log$ of the two sides
are equal.

To compute these logarithmic derivatives, we shall use the formula
\[
  T \frac{d}{dT}\log\det(1-f T)^{-1} = \sum_{n>0} \tr(f^n) T^n \text{,}
\]
whose proof (in a slightly more general setting) is recalled below. On the
left-hand side, one finds
\begin{align*}
  T\frac{d}{dT}\log L(X_0,\sF_0)
    &= \sum_{x\in |X_0|} \sum_n [k(x):\dF_q] \tr\left(F_x^n,\sF\right) T^{n\cdot[k(x):\dF_q]} \\
    &= \sum_{n>0} T^n \sum_{x\in X^{F^n}} \tr\left({F^n}^*,\sF_x\right)
\end{align*}
and on the right-hand side
\[
  \sum_{n>0} T^n \sum (-1)^i \tr\left({F^n}^*,\h_c^i(X,\sF)\right) \text{ ; }
\]
the terms are compared one by one.






\begin{proposition_}\label{II:3-3}
Let $A$ be a commutative ring, $M$ a projective $A$-module of finite type,
and $f$ an endomorphism of $M$. For every invertible formal series $Q$, set
$\frac{d}{dT}\log Q = Q^{-1}\frac{d}{dT}Q$. One has
\begin{equation*}\tag{3.3.1}\label{II:eq:3-3-1}
  T\frac{d}{dT}\log\det(1-f T)^{-1} = \sum_{n>0} \tr(f^n) T^n \text{.}
\end{equation*}
\end{proposition_}

Let $M'$ be an $A$-module such that $M\oplus M'$ is free of finite rank.
Replacing $(M,f)$ by $(M\oplus M',f\oplus 0)$ changes neither side of
(\ref{II:3-3}), and reduces us to the case where $M$ is free. For $M=A^d$,
\eqref{II:eq:3-3-1} is an algebraic identity in the coefficients $f_i^j$ of
$f$. By the principle of extension of algebraic identities, it suffices to
verify it for $A$ an algebraically closed field of characteristic $0$. Since
both sides are additive in $M$ in an extension
$0\to M'\to M\to M''\to 0$, it even suffices to treat only the case where
$M$ has rank $1$. If $f$ is multiplication by $a$,
\eqref{II:eq:3-3-1} then reduces to the formula
\[
  T\frac{d}{dT}\log\left((1-a T)^{-1}\right) = \sum_{n>0} a^n T^n \text{.}
\]






\begin{corollary_}\label{II:3-4}
Let $n$ be an integer, and let $f^{(n)}$ be the endomorphism
$(x_1,\dotsc,x_n)\mapsto \left(f(x_n),x_1,\dotsc,x_{n-1}\right)$ of $M^n$.
One has
\begin{equation*}\tag{3.4.1}\label{II:eq:3-4-1}
  \det(1-f T^n) = \det(1-f^{(n)} T) \text{.}
\end{equation*}
\end{corollary_}

Proceeding as in (\ref{II:3-3}), we reduce to assuming that $A$ has
characteristic $0$. It then suffices to verify that the logarithmic
derivatives $-T\frac{d}{dT}\log$ of the two sides are equal.
\begin{align*}
  -T\frac{d}{dT}\log\det(1-f T^n) &= - n T^n \frac{d}{dT^n} \log\det(1-f T^n) = n \sum_{m>0} \tr(f^m) T^{n m} \text{,} \\
  -T \frac{d}{dT}\log\det\left(1-f^{(n)} T\right) &= \sum_{m>0} \tr\left(f^{(n) m}\right) T^m \text{,}
\end{align*}
and one observes that
\begin{align*}
  \tr(f^{(n) m}) &= 0 \qquad \text{if $n\nmid m$} \\
  \tr(f^{(n) n m}) &= n \tr(f^m) \text{.}
\end{align*}






\subsection{Remark}\label{II:3-5}

For $A=\dQ_\ell$, \eqref{II:eq:3-4-1} is the special case
$X_0=\spec(\dF_{q^n})$ of (\ref{II:3-1}). Using this identity, it is easy to
verify directly that the right-hand side of (\ref{II:3-1}) is independent of
$q$ and of the structure morphism $X_0\to\spec(\dF_q)$.






\subsection{Remark}\label{II:3-6}

If $j:X_0\hookrightarrow \bar X_0$ is a compactification of $X_0$, one has
$L(X_0,\sF_0) = L(\bar X_0, j_!\sF_0)$ and
$\h_c^i(X,\sF) = \h^i(\bar X,j_!\sF)$. This reduces (\ref{II:3-1}) to the
proper case, and explains the use of cohomology with proper support. For $X$
proper, (\ref{II:3-2}) has the form of a Lefschetz trace formula. Note that
the local terms $\tr(F^{n *},\sF_x)$ ($x\in X^{F^n}$) depend only on the
fiber of $\sF$ at $x$ and carry no multiplicity. For $X$ and $\sF$ smooth,
this corresponds to the fact that the graph of $F$ is transverse to the
diagonal (since $d F = 0$).






\subsection{Remark}\label{II:3-7}

If the formalism of $\dQ_\ell$-sheaves is admitted (cf. \ref{II:2-11}; one
needs the Leray spectral sequence and the base-change theorem for the higher
direct images $\eR^q f_!$), it is easy to reduce the proof of
(\ref{II:3-1}, \ref{II:3-2}) to the case where $X_0$ is a smooth curve and
$\sF_0$ is lisse. This is clearly explained in \cite[\S 5]{gr95} (for
\ref{II:3-1}; \ref{II:3-2} is treated in the same way).






\subsection{}\label{II:3-8}

There are two methods for proving (\ref{II:3-2}).



\paragraph{Lefschetz--Verdier}
If $X_0$ is proper (cf. \ref{II:3-6}), the general Lefschetz--Verdier trace
formula expresses the right-hand side of (\ref{II:3-2}) as a sum of local
terms, one for each point of $X^{F^n}$. In the original version of
\cite{sga5}, this formula was proved only modulo resolution of singularities.
The reader will find an unconditional proof in the definitive version. In
the case of curves, to which one can reduce (\ref{II:3-7}), all the
ingredients of the proof were in fact already available.

To deduce (\ref{II:3-2}) from the Lefschetz--Verdier formula, one must be
able to compute its local terms. For a curve and the Frobenius endomorphism,
this had been done by Artin and Verdier \cite{ve67}.



\paragraph{Nielson--Wecken}
A method inspired by the work of Nielsen and Wecken makes it possible to
reduce (\ref{II:3-2}) to a special case proved by Weil; this is what will be
explained in the following sections.











\section{Reduction to theorems in \texorpdfstring{$\dZ/\ell^n$}{Z/l n}-cohomology}\label{II:4}

We shall prove (\ref{II:3-2}) by passage to the limit from an analogous
theorem in $\dZ/\ell^n$-cohomology. Its statement presents the following
difficulty. Let $X_0$ be a separated scheme of finite type over $\dF_q$, and
let $\sF_0$ be a flat constructible sheaf of $\dZ/\ell^n$-modules. Its fibers
are finite-rank free $\dZ/\ell^n$-modules, and the left-hand side of
(\ref{II:3-2}) therefore makes sense (it is an element of $\dZ/\ell^n$).
The $\h_c^i(X,\sF)$, on the other hand, will not in general be free, so the
right-hand side is not defined. Overcoming this difficulty is one of
Grothendieck's essential uses of the theory of derived categories.






\subsection{}\label{II:4-1}

For the definition of the derived category $\eD(\cA)$ of an abelian category
$\cA$, and that of the subcategories $\eD^+(\cA)$, $\eD^-(\cA)$,
$\eD^b(\cA)$, I refer to Verdier's text in this volume. For $\Lambda$ a ring,
write $\eD(\Lambda)$ instead of
$\eD(\text{category of left $\Lambda$-modules})$; for $X$ a site, write
$\eD(X,\Lambda)$ instead of
$\eD($category of sheaves of left $\Lambda$-modules$)$. Likewise for
$\eD^-$, etc\ldots

The sign $\eR$ (resp. $\mathsf{L}$) indicates a right-derived (resp.
left-derived) functor. For example, $\lotimes$ denotes the derived functor of
the tensor product: for $K$ in $\eD^-(\Lambda^\circ)$ and $L$ in
$\eD^-(\Lambda)$ ($\Lambda^\circ$ the opposite ring of $\Lambda$),
$K\lotimes L$ is computed as follows: take quasi-isomorphisms $K'\to K$,
$L'\to L$, with $K'$ and $L$ bounded above and one of them having flat
components, and take the simple complex $s(K'\otimes_\Lambda L)$ associated
with the double complex $K'\otimes_\Lambda L'$.






\subsection{}\label{II:4-2}

Even to treat only the case of $\dZ/\ell^n$-cohomology, we shall need the
theory of traces of endomorphisms of modules over rings that are not
necessarily commutative (finite-group algebras $\dZ/\ell^n[G]$). These
traces were introduced by Stallings and Hattori \cite{ba76}.

Let $\Lambda$ be a ring (with identity, but not necessarily commutative).
Denote by $\Lambda^\natural$ the quotient of the additive group of $\Lambda$
by the subgroup generated by the commutators $a b-b a$. If $f$ is an
endomorphism of the free left $\Lambda$-module $\Lambda^n$, with matrix
$f_i^j$, denote by $\tr(f)$ the image of $\sum_i f_i^i$ in
$\Lambda^\natural$. If $f$ and $g$ are two homomorphisms
$f:\Lambda^n \leftrightarrows \Lambda^m : g$, one immediately verifies that
$\tr(f g) = \tr(g f)$.

Let $f$ be an endomorphism of a projective $\Lambda$-module of finite type
$P$. Write $P$ as a direct factor of a free module $\Lambda^n$: this means
choosing a module $P'$ and an isomorphism
$\alpha:P\oplus P'\iso \Lambda^n$, or equivalently choosing a homomorphism
$a:P\to \Lambda^n$ and a retraction $b:\Lambda^n\to P$ (take
$P'=\ker(b)$). Let $f'=\alpha(f\oplus 0)\alpha^{-1} =af b$. Set
$\tr(f) = \tr(f')$. This element of $\Lambda^\natural$ depends only on $f$:
if another choice $c:P\leftrightarrows \Lambda^m:d$ is made, one has
$a=a(dc)(ba)$, whence
\[
 \tr(a f b) = \tr(adcbafb) = \tr(cbafbad) = \tr(cfd) \text{.}
\]

If $f$ is an endomorphism of a $\dZ/2$-graded projective
$\Lambda$-module of finite type, with components $f_i^j:P^j\to P^i$, set
\[
  \tr(f) = \tr(f_0^0) - \tr(f_1^1) \text{.}
\]
If there is a risk of ambiguity, we shall instead write $\tr(f;P)$. If the
homomorphisms $f:P\leftrightarrows Q:g$ are homogeneous, reduction to the
free case immediately verifies that
\begin{equation*}\tag{4.2.1}\label{II:eq:4-2-1}
  \tr(f g) = \tr\left((-1)^{\deg f \deg g} g f\right) \text{.}
\end{equation*}

If $f$ is an endomorphism of a filtered $\dZ/2$-graded $\Lambda$-module
$(P,F)$, with finite filtration compatible with the grading, and the
$\gr_F^i(P)$ are projective of finite type, then $P$ is projective of finite
type, and
\begin{equation*}\tag{4.2.2}\label{II:eq:4-2-2}
  \tr(f;P) = \sum \tr\left(f; \gr_F^i(P) \right)
\end{equation*}
(split the filtration to reduce to the case of a direct sum).






\subsection{}\label{II:4-3}

A $\dZ$-graded module will be regarded as $\dZ/2$-graded by the parity of
the degree. In particular, if $f$ is an endomorphism of a bounded complex of
projective $\Lambda$-modules of finite type, set
\begin{equation*}\tag{4.3.1}\label{II:eq:4-3-1}
  \tr(f) = \sum (-1)^i \tr(f^i) \text{.}
\end{equation*}

By \eqref{II:eq:4-2-1}, if $f$ is homotopic to zero, $\tr(f) = 0$:
\begin{equation*}\tag{4.3.2}\label{II:eq:4-3-2}
  f = d H + H d \rightarrow \tr(f) = 0 \text{.}
\end{equation*}

Let $\mathsf{K}_\text{parf}(\Lambda)$ be the category whose objects are
bounded complexes of projective $\Lambda$-modules of finite type and whose
arrows are morphisms of complexes taken up to homotopy. The functor
$\mathsf{K}_\text{parf}(\Lambda)\to\eD(\Lambda)$ is fully faithful. Denote its
essential image by $\mathsf{D}_\text{parf}(\Lambda)$. Formulas
\eqref{II:eq:4-3-1} and \eqref{II:eq:4-3-2} make it possible to define
$\tr(f)$ for $f$ an endomorphism of
$K\in\ob \mathsf{D}_\text{parf}(\Lambda)$ (and this trace depends only on the
isomorphism class of $(K,f)$).






\subsection{}\label{II:4-4}

Recall that the filtered derived category $\mathsf{DF}(\cA)$ of an abelian
category $\cA$ is the category obtained from that of complexes of objects of
$\cA$ with finite filtration by inverting the filtered quasi-isomorphisms
($=$ morphisms of filtered complexes inducing a quasi-isomorphism on the
associated graded). There are functors ``underlying complex'':
$\mathsf{DF}(\cA)\to\eD(\cA): (K,F)\mapsto K$, and ``$p$-th component of the
associated graded'':
$\mathsf{DF}(\cA)\to\eD(\cA): (K,F)\mapsto \gr_F^p(K)$.

Let $\mathsf{KF}_\text{parf}(\Lambda)$ be the category whose objects are
bounded filtered complexes $(K,F)$ with finite filtration, with the
$\gr_F^p(K^q)$ projective of finite type, and whose arrows are
filtration-preserving morphisms of complexes taken up to a
filtration-preserving homotopy. The functor
$\mathsf{KF}_\text{parf}(\Lambda)\to\mathsf{DF}(\Lambda)$ is fully faithful;
its essential image $\mathsf{DF}_\text{parf}(\Lambda)$ consists of the $(K,F)$ in
$\mathsf{DF}(\Lambda)$ such that the $\gr_F^p(K)$ belong to
$\mathsf{D}_\text{parf}(\Lambda)$ and are almost all zero; if $(K,F)$ belongs to
$\mathsf{DF}_\text{parf}(\Lambda)$, then $K$ belongs to
$\mathsf{D}_\text{parf}(\Lambda)$.

By \eqref{II:eq:4-2-2}, if $f$ is an endomorphism of
$(K,F)\in\ob\mathsf{DF}_\text{parf}(\Lambda)$, one has
\begin{equation*}\tag{4.4.1}\label{II:eq:4-4-1}
 \tr(f,K) = \sum_p \tr(f,\gr_F^p(K)) \text{.}
\end{equation*}






\subsection{}\label{II:4-5}

Let $\Lambda$ be a ring, and let $K\in\ob\eD^-(\Lambda)$. We say that $K$
has tor-dimension $\leqslant r$ if, for every right $\Lambda$-module $N$, the
hypertors $\h^i(N\lotimes_\Lambda K)$ vanish for $i<-r$. It is said to have
finite tor-dimension if this condition holds for some $r$.

The same terminology is used for a complex of sheaves
$K\in\ob\eD^-(X,\Lambda)$.






\begin{lemma}\label{II:4-5-1}
Let $\Lambda$ be a left Noetherian ring, and let
$K\in\ob\eD^-(\Lambda)$. Then $K$ belongs to
$\mathsf{D}_\text{parf}(\Lambda)$ if and only if it has finite tor-dimension
and the $\h^i(K)$ are of finite type.
\end{lemma}

After replacing $K$ by an isomorphic complex (in $\eD^-(\Lambda)$), we may
suppose that $K$ is bounded above and has finite-rank free components (cf.
\ref{II:4-7} below). If $K$ has tor-dimension $\leqslant r$, the
$\h^i(K)$ vanish for $i<-r$, $K^{-r}/\im(d)$ admits the free resolution
\[
  (\cdots \to K^{-r-1} \to K^{-r}) \to K^{-r}/\im(d) \text{,}
\]
and for every right $\Lambda$-module $N$,
$\h^{-r-k}(N\lotimes_\Lambda K) = \h^{-r-k}(N\otimes_\Lambda K)
= \tor_k(N,K^{-r}/\im d)$ for $k\geqslant 1$.

By hypothesis, these groups vanish; the module $K^{-r}/\im d$ is therefore
flat and finitely presented, i.e. projective of finite type. The quotient
\[
  0 \to K^{-r}/\im(d) \to K^{-r+1} \to K^{-r+2}\to \cdots
\]
of $K$ is quasi-isomorphic to $K$, and this shows that
$K\in\ob\mathsf{D}_\text{parf}(\Lambda)$.






\begin{prop-def_}\label{II:4-6}
Let $X$ be a Noetherian scheme, and let $\Lambda$ be a left Noetherian ring.
Denote by $\eD_\textnormal{ctf}^b(X,\Lambda)$ the subcategory of
$\eD^-(X,\Lambda)$ formed by the complexes $\sK$ satisfying the following
equivalent conditions:
\begin{enumerate}[\indent i)]
  \item $\sK$ is isomorphic (in $\eD^-(X,\Lambda)$) to a bounded complex of
    flat constructible sheaves of $\Lambda$-modules.
  \item $\sK$ has finite tor-dimension and the sheaves $\h^i(K)$ are
    constructible.
\end{enumerate}
\end{prop-def_}

The same argument as in (\ref{II:4-5}), starting from the following lemma.






\begin{lemma_}\label{II:4-7}
If a complex $\sK$ of sheaves of $\Lambda$-modules is such that the
$\h^i(\sK)$ are constructible and vanish for large $i$, there is a
quasi-isomorphism $\sK'\to \sK$, with $\sK'$ bounded above and having flat
constructible components.
\end{lemma_}

Construct $\sK'$ by proceeding from right to left. For an $m$ such that the
$\h^i(\sK)$ ($i\geqslant m$) vanish, take ${\sK'}^i = 0$ for
$i\geqslant m$. Next suppose that the ${\sK'}^i$ have already been
constructed for $i>n$:
\[\xymatrix{
  \cdots \ar[r]
    & \sK^{n-1} \ar[r]
    & \sK^n \ar[r]
    & \sK^{n+1} \ar[r]
    & \cdots \\
  & & & {\sK'}^{n+1} \ar[u] \ar[r]
    & \cdots \text{,}
}\]
with $\h^i(\sK')\iso \h^i(\sK)$ for $i>n+1$, and
$\ker(d)\to\h^{n+1}(\sK)$ surjective in degree $n+1$. Define sheaves $\sA$
and $\sB$ by the Cartesian diagram
\[\xymatrix{
  \sK^n \ar[r]
    & \sK^n/\im(d) \ar[r]
    & \ker(d) \\
  \sA \ar[u] \ar[r]^u
    & \sB \ar[u]\ar[r]
    & \ker(d) \ar[u] \text{.}
}\]

Then $\sB$ is constructible, $u$ is surjective, and to advance one step it
suffices to construct ${\sK'}^n$ flat and constructible and
$v:{\sK'}^n\to \sA$ such that $u v$ is surjective. The possibility of doing so
is guaranteed by the following lemma.






\begin{lemma_}\label{II:4-8}
Let $v:\sA\to \sB$ be an epimorphism of sheaves of $\Lambda$-modules, with
$\sB$ constructible. There is then $u:\sK\to \sA$, with $\sK$ flat and
constructible, such that $v u$ is an epimorphism.
\end{lemma_}

The sheaf $\sA$ is a quotient of a sum of sheaves of the form
$\varphi_! \Lambda$, for $\varphi:U\to X$ étale. Since $\sB$ is Noetherian
(\cite[IX.2.10]{sga4}), a finite sum of these already maps onto $\sB$, which
proves the lemma.






\begin{theorem_}\label{II:4-9}
Let $f:X\to Y$ be a separated morphism of finite type of Noetherian schemes.
If $\sK\in\ob\eD_\textnormal{ctf}^b(X,\Lambda)$, then
$\eR f_!\sK\in \ob\eD_\textnormal{ctf}^b(Y,\Lambda)$.
\end{theorem_}

Recall that if $j:X\hookrightarrow \bar X \xrightarrow{\bar f} Y$ is a
relative compactification of $X/Y$, then by definition
$\eR f_! = \eR \bar f_* \circ j_!$, where $\eR \bar f_*$ is the derived
functor of $\bar f_*$. This formula reduces the proof of (\ref{II:4-9}) to
the case where $f$ is proper. For $f$ proper, $f_*$ has finite cohomological
dimension; this makes it possible to define $\eR f_*$ on the entire derived
category, and also as a functor
\[
  \eR f_* : \eD^-(X,\Lambda) \to \eD^-(Y,\Lambda) \text{.}
\]
By composition, one likewise defines in general
\[
  \eR f_! : \eD^-(X,\Lambda) \to \eD^-(Y,\Lambda) \text{.}
\]

For every right $\Lambda$-module $N$, one has (proof given below)
\begin{equation*}\tag{4.9.1}\label{II:eq:4-9-1}
  N\lotimes \eR f_! \sK \iso \eR f_! (N\lotimes \sK) \text{.}
\end{equation*}

Let us deduce (\ref{II:4-9}) from \eqref{II:eq:4-9-1}.
\begin{enumerate}[\indent a)]
  \item The spectral sequence
    \[
      E_2^{p q} = \eR^p f_! \h^q(\sK) \rightarrow \h^{p+q} \eR f_! \sK \text{.}
    \]
    and the finiteness theorem for $\eR f_!$ show that
    $\eR f_!\sK\in \ob \eD^b(Y,\Lambda)$ and has constructible cohomology.
    It remains to prove finite tor-dimension (\ref{II:4-6}).
  \item If $\sK$ has tor-dimension $\leqslant -r$, the same is true of
    $\eR f_! \sK$: one has
    $\h^i(N\lotimes \eR f_! \sK) = \h^i(\eR f_!(N\lotimes \sK))$, and in the
    spectral sequence above for $N\lotimes \sK$, the $E_2^{p q}$ vanish for
    $q<r$, hence a fortiori for $p+q<r$.
\end{enumerate}

Let us prove \eqref{II:eq:4-9-1}, remaining at the level of complexes.
\begin{enumerate}[\indent a)]
  \item Using the defining formula $\eR f_! = \eR f_*\circ j_!$, we reduce to
    assuming that $f$ is proper.
  \item Represent $\sK$ by a complex bounded above whose components are
    acyclic for $f_*:\eR^p f_* \sK^q = 0$ for $p>0$. One then has
    $\eR f_* K \sim f_* \sK$. It is possible to work in this way with
    complexes bounded above only because $f_*$ has finite cohomological
    dimension.
  \item If $L$ is a free right $\Lambda$-module, one has
    $\eR^p f_*(L\otimes \sK^q) = L\otimes \eR^p f_* \sK^q$: this is trivial
    when $\Lambda$ is free of finite type, and the general case follows by
    passage to the limit. The sheaf $L\otimes \sK^q$ is acyclic for $f_*$,
    and $f_*(L\otimes \sK^q) = L\otimes f_* \sK^q$.
  \item Let $N_\bullet$ be a free resolution of $N$. One has
    $N\lotimes \sK\sim s(N_\bullet\otimes \sK)$ (the simple complex
    associated with the double complex $N_\bullet\otimes \sK$). By c),
    \[
      \eR f_*(N\lotimes \sK) \sim \eR f_* (s(N_\bullet\otimes \sK)) \sim f_*(s(N_\bullet\otimes \sK)) \sim s(N_\bullet\otimes f_* \sK) \sim N\lotimes \eR f_* \sK \text{,}
    \]
    which is the asserted formula.
\end{enumerate}

When $Y$ is the spectrum of a separably closed field, the functor $\eR f_!$
is identified with a functor denoted by
$\eR \Gamma_c : \eD_\text{ctf}^b(X,\Lambda)\to\eD_\text{parf}(\Lambda)$.






\begin{theorem_}\label{II:4-10}
With the notation of \S\ref{II:1}, let $X_0$ be a separated scheme of finite
type over $\dF_q$, and let $\Lambda$ be a Noetherian torsion ring annihilated
by an integer prime to $p$. Let
$\sK_0\in\ob\eD_\textnormal{ctf}^b(X_0,\Lambda)$. With the definition of
traces in (\ref{II:4-3}), one has for every $N$
\begin{equation*}\label{II:eq:4-10}
  \sum_{x\in X^{F^n}} \tr\left({F^n}^*,\sK_x\right) = \tr\left({F^n}^*,\eR\Gamma_c(X,\sK)\right) \text{.}
\end{equation*}
\end{theorem_}

In the remainder of this \S, we shall deduce \ref{II:3-2} from
\ref{II:4-10}, and prove \ref{II:2-10}.






\subsection{Proof of \texorpdfstring{\ref{II:2-10}}{2.10}}\label{II:4-11}

Let $X$ be a separated scheme of finite type over an algebraically closed
field $k$, and let $\sG$ be a $\dQ_\ell$-sheaf on $X$. By \ref{II:2-8}, it
can be written in the form $\sG=\sF\otimes\dQ_\ell$, with $\sF$ a
torsion-free $\dZ_\ell$-sheaf. Set
\[
  K_n = \eR \Gamma_c(X,\sF\otimes\dZ/\ell^{n+1})\in\eD(\dZ/\ell^{n+1}) \text{.}
\]
By \ref{II:4-9}, one has
$K_n\in\ob\eD_\text{parf}(\dZ/\ell^{n+1})$, and it follows from
\eqref{II:eq:4-9-1} that the natural arrow
\begin{equation*}\tag{4.11.1}\label{II:eq:4-11-1}
\xymatrix{
  K_n\lotimes_{\dZ/\ell^{n+1}} \dZ/\ell^n \ar[r] & K_{n-1} \text{,}
}
\end{equation*}
is an isomorphism. The problem of defining this arrow is discussed in
\ref{II:4-12}.

We may now invoke \cite[XV.3.3]{sga5} (p. 31, lines 1--2 from the end; this
passage is independent of the rest of \cite{sga5}) to realize each $K_n$ by
a bounded complex of finite-rank free $\dZ/\ell^{n+1}$-modules, and the
arrows \eqref{II:eq:4-11-1} by \emph{isomorphisms of complexes}
\[\xymatrix{
  K_n\otimes \dZ/\ell^n \ar[r]^-\sim
    & K_{n-1} \text{.}
}\]
Let $K$ be the complex $\varprojlim K_n$. It is a bounded complex of free
$\dZ_\ell$-modules, and
$K_n\simeq K\otimes_{\dZ_\ell}\dZ/\ell^{n+1}$. One has
\[
  \h^i(K) = \varprojlim \h^i(K_n) \text{,}
\]
(indeed, every inverse system of finite groups satisfies the Mittag--Leffler
condition), and $\h^i(K)$ is a $\dZ_\ell$-module of finite type. One has
\[
  \h_c^i(X,\sG) = \left(\varprojlim \h^i(K_n) \right)\otimes_{\dZ_\ell} \dQ_\ell \text{:}
\]
this is a finite-dimensional $\dQ_\ell$-vector space.






\subsection{The arrow \texorpdfstring{\eqref{II:eq:4-11-1}}{(4.11.1)}}\label{II:4-12}

To \emph{define} the arrow \eqref{II:eq:4-11-1}, one cannot proceed as in
\ref{II:4-9}, i.e. resolve $\dZ/\ell^n$ by a complex of free
$\dZ/\ell^{n+1}$-modules, since one wishes to construct an isomorphism in
$\eD_\text{parf}(\dZ/\ell^n)$. The desired arrow is a special case of the
following scalar-extension arrow.

(*) Let $\Lambda\to\Lambda'$ be a homomorphism of Noetherian torsion rings.
For $X$ a separated scheme of finite type over an algebraically closed $k$,
and $\sK\in\ob\eD_\text{ctf}(X,\Lambda)$, one has an isomorphism
\[\xymatrix{
  \eR\Gamma_c(X,\sK)\lotimes_\Lambda \Lambda' \ar[r]^-\sim & \eR \Gamma_c(X,\sK\lotimes_\Lambda \Lambda') \text{.}
}\]

More general arrows are constructed in \cite[XVII 4.2.12]{sga4}. The method
is as follows:
\begin{enumerate}[\indent a)]
  \item The definition of $\eR\Gamma_c$ and the formula
    $j_!(\sK\otimes_\Lambda\Lambda') = (j_! \sK)\otimes_\Lambda\Lambda'$ for
    an open immersion reduce us to the case where $X$ is proper.
  \item For $X$ proper, one must replace $\sK$ by a bounded complex whose
    components are acyclic for the functor $\Gamma$, and whose fiber at every
    geometric point (or merely every closed point) is homotopic to a complex
    of flat $\Lambda$-modules. For such a complex, one has
    $\Gamma(X,\sK)\sim \eR\Gamma(X,\sK)$ and
    $K\otimes_\Lambda\Lambda'\sim \sK\lotimes_\Lambda\Lambda'$; the arrow
    \eqref{II:eq:4-11-1} is defined as the composite
    \[
      \eR\Gamma(X,\sK)\lotimes_\Lambda \Lambda' \to \Gamma(X,\sK)\otimes_\Lambda\Lambda' \to \Gamma(X,\sK\otimes_\Lambda\Lambda') \xleftarrow\sim \eR\Gamma(X,\sK\lotimes_\Lambda\Lambda') \text{.}
    \]
  \item It remains to show that $\sK$ can be replaced by a complex of the
    desired kind \cite[XVII.4.2.10]{sga4}. First replace $\sK$ by a bounded
    complex with flat components (\ref{II:4-6}), then take a truncated
    canonical flasque resolution of it (this does not modify the homotopy
    type of the fiber complexes, and one truncates far enough for the
    acyclicity condition for $\Gamma$ to hold (i.e. beyond the cohomological
    dimension of $\Gamma$)).
\end{enumerate}






\subsection{Proof of \texorpdfstring{(\ref{II:3-2})}{(3.2)}}\label{II:4-13}

To simplify notation, we shall treat only the case $n=1$ of
(\ref{II:3-2}). The general case follows by replacing $F$ below by $F^n$.

Thus let $X_0$ be separated and of finite type over $\dF_q$, and let $\sG_0$
be a $\dQ_\ell$-sheaf on $X_0$. One has
$\sG_0=\sF_0\otimes\dQ_\ell$, for a suitable torsion-free
$\dZ_\ell$-sheaf $\sF_0$. Let
$K_n=\eR\Gamma_c(X,\sF\otimes\dZ/\ell^n)$. The Frobenius morphisms induce
morphisms
\[
  F^* : K_n\to K_n \qquad \text{(in $\eD_\text{parf}(\dZ/\ell^{n+1})$),}
\]
which are deduced from one another via the isomorphisms
\eqref{II:eq:4-11-1}.

As in (\ref{II:4-11}), realize the $K_n$ and the isomorphisms
\eqref{II:eq:4-11-1} at the level of complexes. The already cited passage
\cite[XV.3.3]{sga5} makes it possible to realize the morphisms $F^*$ by
endomorphisms of complexes, again denoted by $F^*$, which reduce to one
another. Again denote by $F^*$ their inverse limit and the endomorphisms
deduced from it. Since
$\h_c^i(X,\sG) = \h^i(K)\otimes\dQ_\ell = \h^i(K\otimes\dQ_\ell)$, one has
(with the sign convention of \ref{II:eq:4-3-1})
\[
  \tr\left(F^*, \h_c^\bullet(X,\sG)\right) = \tr\left(F^*,K^\bullet\otimes\dQ_\ell\right) = \varprojlim_n \tr\left(F^*,K_n^\bullet\right) = \varprojlim_n \tr\left(F^*,\eR\Gamma_c(X,\sF\otimes\dZ/\ell^{n+1})\right) \text{.}
\]
Apply (\ref{II:4-10}) to $\sF_0\otimes\dZ/\ell^{n+1}$ (viewed as a complex
concentrated in degree $0$), for $\Lambda=\dZ/\ell^{n+1}$. One finds that
\begin{align*}
  \tr\left(F^*,\eR\Gamma_c(X,\sF\otimes\dZ/\ell^{n+1})\right)
    &= \sum_{x\in X^F} \tr\left(F^*,\sF_x\otimes\dZ/\ell^{n+1}\right) \\
    &= \sum_{x\in X^F} \tr\left(F^*,\sG_x\right) \mod \ell^{n+1} \text{,}
\end{align*}
and (\ref{II:3-2}) follows by passage to the limit.











\section{The Nielsen--Wecken method}\label{II:5}

In this chapter, we prove two special cases of (\ref{II:4-10}). The general
case will be considered in the next chapter.






\begin{lemma_}\label{II:5-1}
Theorem (\ref{II:4-10}) is true for $X_0$ of dimension $0$.
\end{lemma_}

In dimension $0$, formula \eqref{II:eq:4-10} holds for every correspondence,
and not only for the iterates of a Frobenius. It is a special case of the
following statement (applied to $X(\dF)$, $K$, and the ${F^*}^n$).

(*) Let $X$ be a finite set, $F:X\to X$,
$\sK\in\ob\eD_\text{parf}(X,\Lambda)$, and $F^*:F^* \sK\to \sK$. One has
\[
  \sum_{X^F} \tr\left(F^*,\sK_x\right) = \tr\left(F^*,\Gamma(X,K)\right) \text{.}
\]

One has $\Gamma(X,\sK)=\bigoplus_{x\in X} \sK_x$, and the formula follows
from the fact that for every morphism
$u:\bigoplus_{x\in X} \sK_x\to \bigoplus_{x\in X} \sK_x$, with matrix
$u_x^y$, one has $\tr(u)=\sum_{x\in X} \tr\left(u_x^x\right)$.






\begin{lemma_}\label{II:5-2}
Let $\Lambda$ be a Noetherian torsion ring annihilated by a power of the prime
number $\ell\ne p$, let $X_0$ be a smooth projective connected curve over
$\dF_q$, let $j:U_0\hookrightarrow X_0$ be a dense open subscheme of $X_0$,
and let $\sG_0$ be a locally constant sheaf of projective
$\Lambda$-modules on $U_0$. One has
\[
  \sum_{x\in U^F} \tr\left(F^*,\sG_x\right) = \tr\left(F^*,\eR\Gamma_c(U,\sG)\right) \text{.}
\]
\end{lemma_}

More precisely, we shall deduce (\ref{II:5-2}) from the following theorem.






\begin{theorem_}\label{II:5-3}
Let $X$ be a nonsingular projective curve over an algebraically closed field
$k$, let $f$ be an endomorphism of $X$ with isolated fixed points, and let
$v(f)$ be the number of fixed points of $f$, each counted with its
multiplicity: $v(f)$ is the intersection number of the graph of $f$ with the
diagonal of $X\times X$. Then, for $\ell$ prime to the characteristic
exponent of $k$, one has
\[
  v(f) = \sum (-1)^i \tr\left(f^*,\h^i(X,\dQ_\ell)\right) \text{.}
\]
\end{theorem_}

This theorem is due to Weil (\cite{we48}, n$^\circ$ 43, 66, and 68), and Weil's
proof is reproduced in Lang \cite{la83} (VI \S3 combined with VII \S2,
Th. 3). It can also be deduced from the formalism of the cohomology class
associated with a cycle (\nameref{IV}, \ref{IV:3-7}).






\begin{corollary_}\label{II:5-4}
Let $j:U\hookrightarrow X$ be a dense open subscheme of $X$, such that
$f^{-1}(U)=U$. If the fixed points of $f$ in $X\setminus U$ have multiplicity
one, then the number $v_U(f)$ of fixed points of $f$ in $U$ (each counted
with its multiplicity) is
\[
  v_U(f) = \sum (-1)^i \tr\left(f^*,\h_c^i(U,\dQ_\ell)\right) \text{.}
\]
\end{corollary_}

Let $F$ be the finite set $X\setminus U$. The cohomology exact sequence
defined by the short exact sequence
$0\to j_! \dQ_\ell\to\dQ_\ell\to{\dQ_\ell}_F\to 0$,
\[\xymatrix{
  \ar[r]^-\partial
    & \h_c^i(U,\dQ_\ell) \ar[r]
    & \h^i(X,\dQ_\ell) \ar[r]
    & \h^i(F,\dQ_\ell) \ar[r]^-\partial
    &
}\]
gives, with the sign convention of (\ref{II:3-1}),
\[
  \tr\left(f^*,\h_c^\bullet(U,\dQ_\ell)\right)
    = \tr\left(f^*,\h^\bullet(X,\dQ_\ell)\right) - \tr\left(f^*,\h^\bullet(F,\dQ_\ell)\right)
    = v(f) - \tr(f^*,\dQ_\ell^F) \text{.}
\]
The trace of $f^*$ on $\dQ_\ell^F$ is the number $v_F(f)$ of fixed points of
$f$ in $F$; the multiplicity-one hypothesis therefore ensures that
\[
  \tr\left(f^*,\h_c^\bullet(U,\dQ_\ell)\right)
    = v(f) - v_F(f) = v_U(f) \text{,}
\]
as promised.

The proof of (\ref{II:5-2}) uses some lemmas concerning a special case of
noncommutative traces (\ref{II:4-2}), which we now develop. \emph{For
brevity, we shall say projective for projective of finite type}.






\subsection{}\label{II:5-5}

Let $\Lambda$ be a ring, and let $H$ be a finite group. The map
$\sum a_h \cdot a\mapsto $ class of $a_e$, from the group algebra
$\Lambda[H]$ to $\Lambda^\natural$, vanishes on the commutators $ab-ba$ of
$\Lambda[H]$. It therefore factors through
\[
  \varepsilon : \Lambda[H]^\natural \to \Lambda^\natural \text{.}
\]
If $F$ is an endomorphism of a projective $\Lambda[H]$-module $P$, set
\[
  \tr_\Lambda^H(F) = \varepsilon \tr_{\Lambda[H]} (F)
\]
where the trace of (\ref{II:4-2}) appears on the right-hand side. To avoid
ambiguity, we shall sometimes write this trace as $\tr_\Lambda^H(F,P)$.






\begin{proposition_}\label{II:5-6}
$\tr_\Lambda(F) = |H|\tr_\Lambda^H(F)$.
\end{proposition_}

The formal properties of traces reduce us to treating only the case
$P=\Lambda[H]$. The endomorphism $F$ is then right multiplication by an
element $\sum a_h\cdot h$ of $\Lambda[H]$, and
$\tr_\Lambda(F) = |H| a_e$, $\tr_\Lambda^H(F) = a_e$.






\subsection{}\label{II:5-7}

Let $A$ be a commutative ring, and let $\Lambda$ be an $A$-algebra.
Multiplication by an element of $A$ passes to the quotient to define an
$A$-module structure on $A^\natural$. If $H$ is a finite group, $P$ a
projective $A[H]$-module, and $M$ a $\Lambda[H]$-module projective as a
$\Lambda$-module, it is known that the $\Lambda$-module $P\otimes_A M$,
equipped with the diagonal action of $H$, is a projective
$\Lambda[H]$-module. Indeed:
\begin{enumerate}[\indent a)]
  \item It suffices to verify this for $P=A[H]$.
  \item The $\Lambda$-module $P\otimes_A M$, equipped with the action of $H$
    defined by its action on the first factor, is clearly a projective
    $\Lambda[H]$-module; denote it by $(P\otimes_A M)'$.
  \item The map $h\otimes m\mapsto h\otimes h^{-1} m$ induces an isomorphism
    \begin{equation*}\tag{5.7.1}\label{II:eq:5-7-1}
      A[H]\otimes_A M \iso (A[H]\otimes_A M)' \text{.}
    \end{equation*}
\end{enumerate}






\begin{proposition_}\label{II:5-8}
With the preceding notation, if $u$ is an endomorphism of $P$ and $v$ an
endomorphism of $M$, then
\[
  \tr_\Lambda^H(u\otimes v) = \tr_A^H(u)\tr_\Lambda(v) \text{.}
\]
\end{proposition_}

We reduce to assuming that $P=A[H]$; the endomorphism $u$ is then right
multiplication $m_x$ by an element $x=\sum a_h h$ of $A[H]$. Isomorphism
\eqref{II:eq:5-7-1} takes $u\otimes v$ to
$\sum a_h m_h\otimes h^{-1} v$, whose trace is
$a_e\cdot\tr_\Lambda(v)$.






\subsection{}\label{II:5-9}

Let $G$ be a group that is an extension of $\dZ$ by a finite group $H$
\[
  1 \to H \to G \to \dZ \to 1 \text{,}
\]
and let $G^+$ be the inverse image of $\dN$. Let $\Lambda$ be a ring, and let
$P$ be a $\Lambda[G^+]$-module projective as a $\Lambda[H]$-module. If
$g\in G^+$, denote by $Z_g$ the centralizer of $g$ in $H$; multiplication by
$g$ is an endomorphism of $P$, viewed as a $\Lambda[Z_g]$-module. Since $P$
is $\Lambda[Z_g]$-projective, a trace $\tr_\Lambda^{Z_g}(g)$ is defined. By
\ref{II:5-6}, one has
\begin{equation*}\tag{5.9.1}\label{II:eq:5-9-1}
  \tr_\Lambda(g) = |Z_g| \cdot \tr_\Lambda^{Z_g}(g) \text{.}
\end{equation*}

Suppose that $\Lambda$ is an algebra over a commutative ring $A$. For $P$ an
$A[G^+]$-module projective as an $A[H]$-module, and $M$ a
$\Lambda[G^+]$-module projective as a $\Lambda$-module, the
$\Lambda[G^+]$-module $P\otimes_A M$ (diagonal action of $G^+$) is projective
as a $\Lambda[H]$-module and, by \ref{II:5-8}, for every $g\in G^+$ one has
\begin{equation*}\tag{5.9.2}\label{II:eq:5-9-2}
  \tr_\Lambda^{Z_g}(g,P\otimes_A M) = \tr_A^{Z_g}(g,P) \cdot \tr_\Lambda(g,M) \text{.}
\end{equation*}






\subsection{}\label{II:5-10}

Let $P$ be a $\Lambda[G^+]$-module projective as a $\Lambda[H]$-module, and
let $P_H$ be the coinvariants of $H$ in $P$. This is a projective
$\Lambda$-module. The action of $g\in G^+$ passes to the quotient and defines
an endomorphism of $P_H$ that depends only on the image of $g$ in $\dN$. We
denote by $F$ the action of the elements $g\in G^+$ with image $1$.






\begin{proposition_}\label{II:5-11}
$\tr_\Lambda(F,P_H) = \sum_{g\mapsto 1}' \tr_\Lambda^{Z_g}(g,P)$, where
$\sum_{g\mapsto 1}'$ denotes a sum over the $H$-conjugacy classes of elements
of $G^+$ with image $1$ in $\dN$.
\end{proposition_}

After multiplication by $|H|$, this formula can be verified as follows:
\[
  |H|\cdot \tr_\Lambda(F,P_H)
    = \tr_\Lambda\left(\sum_{g\mapsto 1} g,P\right)
    = \sum_{g\mapsto 1}' \frac{|H|}{|Z_g|}\tr_\Lambda(g,P)
    = \sum_{g\mapsto 1}' |H|\cdot \tr_\Lambda^{Z_g}(g,P) \text{.}
\]

Choose an element $g\in G^+$ with image $1$ in $\dN$, and let $\sigma$ be
the automorphism of $\Lambda[H]$ induced by the automorphism
$\operatorname{ad} g$ of $H$. Giving a $\Lambda[G^+]$-module amounts to
giving a $\Lambda[H]$--module equipped with a $\sigma$-linear endomorphism
$\gamma$: let $g^n h$ act by $\gamma^n h$. In this language, formula
\ref{II:5-11} takes the form: for $P$ a projective $\Lambda[H]$-module and
$\gamma$ a $\sigma$-linear endomorphism of $P$, one has
\[
  \tr_\Lambda(\gamma,P_H) = \sum_h' \tr_\Lambda^{Z_g}(h\gamma,P)\text{,}
\]
where $\sum'$ is a sum over the $\operatorname{ad}g$-conjugacy classes in
$H$ (orbits of $H$ acting on itself by the maps
$x\mapsto h x \operatorname{ad} g(h)^{-1}$).

In this form, the properties of traces immediately reduce us to the case
$P=\Lambda[H]$. The endomorphism $\gamma$ then has the form
$x\mapsto \sigma(x) a$, with $a=\sum a_h h$ in $\Lambda[H]$. It suffices to
treat the universal case where the $a_h$ are indeterminates. In this case,
$\Lambda=\dZ\langle a_h\rangle_{h\in H}$ (noncommutative polynomials), and
$\Lambda^\natural$ is torsion-free: one obtains \ref{II:5-11} by the
argument at the beginning and division by $|H|$. One can also calculate
directly.






\subsubsection{Remark}\label{II:5-11-1}

Everything said above applies equally well to $\dZ/2$-graded modules, or to
complexes of modules (cf. \ref{II:4-2}, \ref{II:4-3}).






\subsection{}\label{II:5-12}

Retain the notation of \ref{II:5-2}, and let $A=\dZ/\ell^n$, with $n$ large
enough that $\Lambda$ is an $A$-algebra. Let $f:V_0\to U_0$ be a finite
\'etale connected covering of $U_0$, Galois with Galois group $H$, and such
that the sheaf $f^*\sG_0$ is constant. Denote its constant value
$\h^0(V_0,f^*\sG_0)$ by $M$; it is a $\Lambda[H]$-module projective as a
$\Lambda$-module.

The sheaf $f_* A$ on $U_0$, equipped with the natural action of $H$, is a
locally free sheaf (of rank one) of $A[H]$-modules. The complex
$\eR\Gamma_c(U,f_* A)$ is therefore an object of
$\eD_\text{parf}\left(A[H]\right)$. It is equipped with a Frobenius
endomorphism $F^*$.

For the natural action of $H$ on $f_* f^*\sG_0$, the trace morphism
$f_* f^*\sG_0\to\sG_0$ factors through an isomorphism
\[
  \left(f_* f^*\sG_0\right)_H \to \sG_0 \text{.}
\]
Moreover, one has $f_* f^*\sG_0=f_* A\otimes_A M$ (with diagonal action of
$H$). The isomorphism
\[\xymatrix{
  \sG_0
    & \ar[l]_-\sim (f_* A\otimes_A M)_H = (f_* A\otimes_A M)\otimes_{\Lambda[H]} \Lambda
}\]
derives (\ref{II:4-12}) to an isomorphism
\[
  \eR\Gamma_c(U,\sG) \sim \left(\eR\Gamma_c(U,f_* A)\otimes_A M\right)\otimes_{\Lambda[H]} \Lambda \text{.}
\]

The endomorphism $F^*$ of the left-hand side is deduced from the
endomorphism $F^*$ of $\eR\Gamma_c(U,f_* A)$. The complex of
$\Lambda[H]$-modules $\eR\Gamma_c(U,f_* A)$, equipped with $F^*$, can be
viewed as a complex of $\Lambda[G^+]$-modules, for $G^+=H\times \dN$.
Formulas \ref{II:5-10} and \eqref{II:eq:5-9-2}, amplified by
\ref{II:5-11-1}, therefore give
\[
  \tr\left(F^*,\eR\Gamma_c(U,\sG)\right) = \sum_h' \tr_A^{Z_h}\left(h F^*,\eR\Gamma_c(U,f_* A)\right)\cdot \tr_\Lambda(h,M)
\]
where $\sum'$ is a sum over the conjugacy classes in $H$. Moreover,
\[
  |Z_h|\cdot \tr_A^{Z_h}\left(h F^*,\eR\Gamma_c(U,f_* A)\right) = \tr_A\left((F h^{-1})^*,\eR\Gamma_c(V,A)\right) \text{.}
\]
This formula remains true for $A=\dZ/\ell^m$, with $m$ increasingly large.
Passing to the limit, one finds that
$\tr_A^{Z_h}\left(h F^*,\eR\Gamma_c(U,f_* A)\right)$ is equal to
\[
  \frac{1}{|Z_h|} \sum (-1)^i \tr\left((F h^{-1})^*,\h_c^i(V,\dQ_\ell)\right)
\]
(an element of $\dZ_\ell$), and
\begin{equation*}\tag{5.12.1}\label{II:eq:5-12-1}
\tr\left(F^*,\eR\Gamma_c(U,\sG)\right) = \sum_h' \frac{1}{|Z_h|} \tr\left((F h^{-1})^*,\h_c^\bullet(V,\dQ_\ell)\right)\cdot \tr_\Lambda(h,M) \text{.}
\end{equation*}






\subsection{}\label{II:5-13}

It remains to apply \ref{II:5-4} (note that if $\bar V$ is the nonsingular
projective completion of $V$, all the fixed points of $F h^{-1}$ have
multiplicity one). This can be done without any calculation:
\begin{enumerate}[\indent A.]
  \item If $U^F=\varnothing$,
    $\tr\left(F^*,\eR\Gamma_c(U,\sG)\right) = 0$. \end{enumerate}

Indeed, $F h^{-1}$ has no fixed point on $V$; \ref{II:5-4} shows that
$\sum (-1)^i \tr\left((F h^{-1})^*,\h_c^i(V,\dQ_\ell)\right)=0$, and one
applies \eqref{II:eq:5-12-1}.
\begin{enumerate}[\indent A.]
\setcounter{enumi}{1}
  \item In the general case, let $j:U'=U\setminus U^F\hookrightarrow U$.
    The filtration $j_! j^* \sG_0\subset \sG_0$ of $\sG_0$ induces a
    filtration of $\eR\Gamma_c(U,\sG)$ with successive quotients
    $\eR\Gamma_c(U',\sG)$ and $\eR\Gamma_c(U^F,\sG)$. By
    \eqref{II:eq:4-4-1} and \ref{II:5-1}, one has
    \[
      \tr\left(F^*,\eR\Gamma_c(U,\sG)\right) = \tr\left(F^*,\eR\Gamma_c(U',\sG)\right) + \sum_{x\in U^F} \tr(F^*,\sG_x) \text{,}
    \]
    and, since ${U'}^F=0$, the first term on the right-hand side is $0$.
\end{enumerate}











\section{End of the proof of \ref{II:4-10}}\label{II:6}

Let us prove \ref{II:4-10} under the additional hypotheses that the ring
$\Lambda$ is annihilated by a power of a prime number $\ell\ne p$, and that
$n=1$. The general case will follow: decompose $\Lambda$ into its
$\ell$-primary components (with $\ell$ running through the prime numbers),
and replace $\dF_q$ by $\dF_{q^n}$ and $X_0$ by
$X_0\otimes_{\dF_q}\dF_{q^n}$.

For $X_0$ separated and of finite type over $\dF_q$, and
$\sK_0\in\ob\eD_\text{ctf}(X_0,\Lambda)$, set
\[
  T'(X_0,\sK_0) = \sum_{x\in X^F} \tr(F^*,\sK_x)
\]
and
\[
  T''(X_0,\sK_0) = \tr(F^*,\eR\Gamma_c(X,\sK)) \text{.}
\]
We must prove the identity
\[
  (1;X_0,\sK_0) \qquad T'(X_0,\sK_0) = T''(X_0,\sK_0) \text{.}
\]
We begin by observing that $T'$ and $T''$ (denoted generically by $T$) obey
the same formalism:

(A) If $j:U_0\hookrightarrow X_0$ is an open subscheme of $X_0$, with
closed complement $Y_0$, one has
\[
  T(X_0,\sK_0) = T(U_0,\sK_0|U_0) + T(Y_0,\sK_0|Y_0) \text{.}
\]
This is clear for $T'$. For $T''$, observe that the filtration
$0\subset j_! j^* \sK_0\subset \sK_0$ of $\sK_0$ induces a filtration of
$\eR\Gamma_c(X,\sK)$ with successive quotients $\eR\Gamma_c(U,\sK)$ and
$\eR\Gamma_c(Y,\sK)$, and invoke \eqref{II:eq:4-4-1} (more precisely stated:
$\sK_0$, equipped with the said filtration, defines an object of
$\mathsf{DF}(X,\Lambda)$; applying $\eR\Gamma_c$ gives an object of
$\mathsf{DF}_\text{parf}(\Lambda)$ whose underlying object is
$\eR\Gamma_c(X,\sK)$ and whose associated graded is as stated).

(B) If $\sK_0$ is a bounded complex with constructible flat components,
\[
  T(X_0,\sK_0) = \sum (-1)^i T(X_0,\sK^i) \text{.}
\]
This is clear for $T'$; for $T''$, apply the argument above to the stupid
filtration of $\sK_0$ to obtain
\[
  T''(X_0,\sK_0) = \sum T''(X_0,\sK_0^i[-i]) = \sum (-1)^i T''(X_0,\sK_0^i) \text{.}
\]

(C) For a morphism $f:X_0\to Y_0$, one has
\[
  T''(X_0,\sK_0) = T''(Y_0,\eR f_! \sK_0)
\]
and the same identity holds for $T'$ if the fibers of $f$ at the rational
points of $Y_0$ satisfy
$\left(1;f^{-1}(y),\sK_0|f^{-1}(y)\right)$: use the fact that
\[
  T'(X_0,\sK_0) = \sum_{y\in Y^F} T'\left(f^{-1}(y), \sK_0|f^{-1}(y)\right)
\]
and
\[
  \eR\Gamma_c(Y, \eR f_! \sK) = \eR\Gamma_c(X,\sK) \text{.}
\]

Using this formalism and \ref{II:4-7}, the proof of \ref{II:4-10} is reduced
to the two special cases \ref{II:5-1} and \ref{II:5-2}.
